\section{Financial evidence: point-in-time portfolio-library retention}
\label{sec:financial-case-studies}
\label{sec:financial-compression-audit}

The financial study asks whether innovation-safe compression can reduce the
active validation burden of portfolio-strategy libraries while preserving a
larger menu of admissible future research options than a budget-matched
frontier-only library. It is a retrospective development-facing panel, not an
independent confirmatory market test. Its contribution is the combination of
point-in-time source construction, an explicit portfolio grammar, protected
proposal decisions, exact structural rechecks, and a transparent held-out null
result.

\subsection{Sequential information boundary}

At each annual origin \(y\), early historical blocks construct the source
portfolio-strategy library and its operating frontier; the subsequent
predecision block estimates compression profiles and solves the safe and
frontier-only arms; year \(y+1\) enumerates closure-enabled proposal options and
freezes both policy choices; and year \(y+2\) supplies the held-out portfolio
return paths. Proposal outcomes do not enter compression. Evaluation outcomes
do not enter compression, option generation, uncertainty estimation, adoption,
or any policy choice. The full candidate ledger retains unavailable and
rejected candidates rather than conditioning the analysis on successful
selection.

The primary universe is formed from the most liquid eligible US common
equities at each compression origin. A separate plain-ETF universe is the
registered replication when its point-in-time eligibility gate passes.
Eligibility never requires survival into the proposal or evaluation window;
observed delisting returns are retained and a terminated position subsequently
moves to cash. The single failed ETF universe gate remains in the registered
denominator. Raw security identifiers, prices, and returns remain inside the
licensed boundary.

Each source library applies a fixed factorial grammar of portfolio-level
signals, filters, holding horizons, allocation rules, gross-exposure targets,
and exits to the same contemporaneous universe. Closure consists of
\AORPointInTimeCapabilities\ declared rule, interface, data, corporate-action,
portfolio-aggregation, and cost/capacity capabilities. These declarations are
more operationally concrete than free-form module tags, but they remain a
modeled grammar: they do not prove that tacit knowledge or every institutional
artifact is frictionlessly recoverable.

\subsection{Compression and protected proposal policy}

The innovation-safe arm minimizes the registered validation-computation burden
subject to exact source-frontier and full declared-closure equality. Its
frontier-only comparator is augmented under the safe arm's realized burden cap
using a frozen preproposal order. The comparator action set is a subset of the
safe action set in every gate-passed cell. Cash is always feasible, and the
safe policy begins with the exact comparator choice. It can replace that choice
only when a simultaneous moving-block lower confidence bound for incremental
proposal value exceeds 0.0025 in annual certainty-equivalent return. The policy therefore
treats closure as an option, not a forced trade.

Across the gate-passed origins, \AORPointInTimePreprocessingSolutions\ safe
solutions are resolved by exact preprocessing and
\AORPointInTimeSolverSolutions\ are returned with HiGHS
\texttt{OPTIMAL} status. Every result is independently reconstructed and
rechecked for mandatory inactive retention, exact table frontier equality,
exact identity-closure equality, and exact integer burden. Those rechecks
certify feasibility and burden; for the solver-resolved cells, global
optimality remains mixed-integer solver evidence rather than a formal or
exhaustive proof.

Table~\ref{tab:point-in-time-financial-study} gives the aggregate structural
and held-out results. Innovation-safe retention reduces the mean source burden
index by \AORPointInTimeEquityBurdenReduction\ in common equities and
\AORPointInTimeETFBurdenReduction\ in ETFs while preserving all registered
capabilities. Relative to the budget-matched frontier-only library, it exposes
on average \AORPointInTimeEquityAddedOptions\ additional common-equity proposal
actions and \AORPointInTimeETFAddedOptions\ additional ETF actions.

\AORPointInTimeFinancialTable

\subsection{Held-out result: options are not exercises}

The protected policy makes \AORPointInTimeEquityPolicyReplacements\
common-equity replacements. Consequently, every common-equity held-out
innovation-safe-minus-frontier-only contrast is exactly zero because the two
policies choose the same action. This is not evidence that the two libraries
are economically equivalent; it is evidence that none of the additional
options clears the proposal-time uncertainty rule.

The ETF replication makes \AORPointInTimeETFPolicyReplacements\ replacement.
That replacement realizes a held-out certainty-equivalent difference of
\AORPointInTimeETFMinimumCE, producing a mean of
\AORPointInTimeETFMeanCE\ over the complete ETF origins. All other complete ETF
pairs choose the same action. The negative observation is retained and is not
offset by reporting a favorable post hoc selector. Full origin-level results,
the prespecified robustness grid, negative-control policies, search-bias
diagnostics, and capacity records appear in Online Resource~1.

The contrast is a postdecision portfolio certainty-equivalent diagnostic under
the registered risk and cost rule. It is not an alpha estimate, a forecast, or
a deployable-performance claim. In particular, the many structural options do
not count as independent financial observations, and the safe policy's
abstention cannot be reinterpreted as positive performance. \subsection{Robustness, licensed-data boundary, and managerial implication}

The proposal robustness analysis retains the full prespecified grid, including
unavailable cells. Common-equity safe and comparator choices remain identical
throughout that grid. ETF differences remain confined to the same origin and
disappear at higher transaction-cost settings. A closure-sham analysis that
was not frozen before proposal outcomes became known is explicitly unavailable,
not reconstructed post hoc. Capacity analysis scores frozen choices without
reselection and retains missing-source, incomplete-path, and hard participation-
cap failures in its denominators.

Licensed daily panels were used locally to construct the point-in-time
universes and portfolio paths. The public artifact redistributes only aggregate
structural manifests, cryptographic seals, permitted summaries, and synthetic
records. It contains no raw CRSP/WRDS rows, security identifiers, or return
paths. Independent replication of the licensed stages therefore requires
separate data access, whereas all published table values can be regenerated
from the committed public-safe manifests.

For research governance, the result is useful precisely because it is not a
success story about returns. A manager can remove substantial declared
validation burden while exactly preserving current opportunities and a much
larger research-option menu. Whether any added option should be exercised is a
separate statistical decision, and a protected selector can rationally abstain.
Removal from the active maintenance set also does not authorize deletion from
legal, compliance, or reproducibility archives. A stronger empirical claim
would require a newly locked market vintage or independent panel, documented
real capability inventories, and monetary maintenance costs calibrated without
using postdecision outcomes.

\subsection{Calibrated closure-option mechanism experiment}
\label{sec:closure-option-mechanism}

The theorem is universal under its assumptions, but it does not say whether a
finite-data policy can recognize a valuable bridge. We therefore isolate that
question in a design-locked synthetic experiment. The innovation-safe and
frontier-only libraries have identical current frontiers and identical
retention burdens. Only the safe library carries the capability needed to run
one declared research project. For delay (d), survival probability \(\rho\),
admission probability \(\pi\), descendant gain \(G\), and research cost
\(\kappa\), the registered treatment margin is
\[
 M=\beta^d\rho^d\pi G-\kappa.
\]
The safe oracle researches exactly when \(M>0\); the frontier-only project is
infeasible. A learner observes an independent noisy proposal signal and adopts
only when it exceeds a threshold calibrated under the null.

The noise law uses only aggregate nuisance estimates from
\AORPointInTimeCompletePaths\ complete proposal paths and
\AORPointInTimeResidualObservations\ within-cell residual observations in the
point-in-time financial study. Cross-strategy date means are removed before
estimating volatility, lag-one dependence, and tail thickness. These
quantities calibrate noise and dependence, not the sign of any regime. The
test-world seeds are disjoint from calibration, and the powered margin is
fixed at twice the null threshold before test simulation. Thus its size is a
power-design dose, not an estimated return premium.

Table~\ref{tab:closure-option-mechanism} reports all four prespecified regimes.
At the low positive dose, the oracle value is positive but the learned-value
interval crosses zero and adoption is only \AORClosureLowDoseAdoption. At the
powered dose \(M=\AORClosurePoweredMargin\), the learned mean is
\AORClosurePoweredLearnedMean\ and adoption is
\AORClosurePoweredAdoption. The null point adoption rate,
\AORClosureNullAdoption, satisfies the registered
\AORClosureNullAdoptionGate\ finite-test gate;
its Wilson upper endpoint is \AORClosureNullWilsonUpper, so the interval is
not wholly below that threshold. Every world passes the equal-frontier,
equal-burden, and bridge-availability structural checks.

\AORClosureOptionTable

This evidence makes the option/exercise distinction concrete. Exact closure
preservation can create access to a genuinely positive project while a
conservative learner rationally fails to exercise a small option at the
available signal-to-noise ratio. Conversely, the powered row shows recovery
only for a deliberately larger treatment. The experiment validates the
declared bridge mechanism under calibrated synthetic noise. It is not a
historical market test, a causal return estimate, an alpha estimate, or a
claim that the powered margin is empirically prevalent.
